\section{Introduction}

Ballot-polling RLAs are often described by their stopping rule, but the public
record also has to support replay. BRAVO, documented in V.12, is per-draw:
each sampled ballot updates a sequential likelihood ratio. Minerva moves the
same ballot-polling idea to audit rounds. The replay question is therefore not
only whether a final stop was declared, but whether each declared round boundary
and round statistic follows from the reported totals, risk limit, and observed
ballots.

RCOUNT treats this as a transcript contract. The verifier does not accept a
vendor-exported risk value as authoritative. It partitions sample steps into
rounds, recomputes the Minerva binomial-tail ratio with exact rational
arithmetic in \texttt{rcount-stats}, and compares the computed round statistic,
p-value in parts per million, and stop decision to the declared audit record.

This paper covers that Minerva slice and the related Athena method surface.
It does not claim that Rhode Island Rep. 28 or any other public audit has
enough released ballot evidence to recompute its Minerva risk value. It claims
that synthetic two-candidate Minerva records with sufficient observations are
replayable, and that Athena records are preserved as method-specific boundaries
rather than silently treated as generic unsupported methods.
